% Filename: slide03.tex
% This code is part of 'Cursos CAMECC: Introducao ao LaTeX para o Curso 29 - Licenciatura em Matematica'
% 
% Description: This file correspond to the slides of using for cap03@camecc_lic.tex of the course.
% 
% Created: 07.06.12 11:30:10 AM
% Last Change: 07.06.12 11:30:30 AM
% 
% Authors:
% - Raniere Silva, r.gaia.cs@gmail.com
% 
% Organization: CAMECC - Centro Academico dos Estudantes do IMECC
% 
% Copyright (c) 2012, Raniere Silva. All rights reserved.
% 
% This work is licensed under the Creative Commons Attribution-ShareAlike 3.0 Unported License. To view a copy of this license, visit http://creativecommons.org/licenses/by-sa/3.0/ or send a letter to Creative Commons, 444 Castro Street, Suite 900, Mountain View, California, 94041, USA.
%
% This work is distributed in the hope that it will be useful, but WITHOUT ANY WARRANTY; without even the implied warranty of MERCHANTABILITY or FITNESS FOR A PARTICULAR PURPOSE.
%
\documentclass[12pt]{beamer}
\usetheme{CambridgeUS}
\let\Tiny=\tiny  % Redefine at least \Tiny for avoid warning
% Filename: use_pack.tex
% This code is part of 'Cursos CAMECC: Introducao ao LaTeX para o Curso 29 - Licenciatura em Matematica'
% 
% Description: This file is only a list of the packages used in the course.
% 
% Created: 07.06.12 11:30:10 AM
% Last Change: 07.06.12 11:30:30 AM
% 
% Authors:
% - Raniere Silva, r.gaia.cs@gmail.com
% 
% Organization: CAMECC - Centro Academico dos Estudantes do IMECC
% 
% Copyright (c) 2012, Raniere Silva. All rights reserved.
% 
% This work is licensed under the Creative Commons Attribution-ShareAlike 3.0 Unported License. To view a copy of this license, visit http://creativecommons.org/licenses/by-sa/3.0/ or send a letter to Creative Commons, 444 Castro Street, Suite 900, Mountain View, California, 94041, USA.
%
% This work is distributed in the hope that it will be useful, but WITHOUT ANY WARRANTY; without even the implied warranty of MERCHANTABILITY or FITNESS FOR A PARTICULAR PURPOSE.
%
\usepackage[brazil]{babel}
\usepackage[utf8]{inputenc}
\usepackage[T1]{fontenc} 
\usepackage{pb-diagram}
\usepackage{graphicx, color}
% \usepackage[top=3cm,left=2cm,right=2cm,bottom=3cm]{geometry} % Must be repeat in the file.
\usepackage{subfig}
\usepackage{enumerate}
\usepackage{algorithmic}
\usepackage{algorithm}
\usepackage{listings}
\usepackage{listingsutf8}
\usepackage{indentfirst}
\usepackage{multirow}
\usepackage{amsthm}
\usepackage{amsmath}
\usepackage{amsfonts}
\usepackage{amssymb}
\usepackage{makeidx}  % For index.
\makeindex  % For index.
\usepackage{tikz}
\usetikzlibrary{patterns}
\usepackage{epsfig}
\usepackage{latexsym}
\usepackage{makeidx}
\usepackage{url}
\usepackage{hyperref}
\usepackage{breakurl}
\usepackage{multicol}
\usepackage{parcolumns}
\usepackage[official]{eurosym}
\lstset{
breaklines=true,
literate={é}{{\'e}}1 {á}{{\'a}}1 {ã}{{\~a}}1,
}

\renewcommand{\lstlistingname}{Código}
\lstset{
language=TeX,                   % the language of the code
basicstyle=\ttfamily\footnotesize,     % the size of the fonts that are used for the code
%numbers=left,                   % where to put the line-numbers
%numberstyle=\footnotesize,      % the size of the fonts that are used for the line-numbers
%stepnumber=5,                   % the step between two line-numbers. If it's 1, each line 
%numbersep=5pt,                  % how far the line-numbers are from the code
showspaces=false,               % show spaces adding particular underscores
showstringspaces=false,         % underline spaces within strings
showtabs=false,                 % show tabs within strings adding particular underscores
tabsize=2,                      % sets default tabsize to 2 spaces
%captionpos=t,                   % sets the caption-position to bottom
breaklines=true,                % sets automatic line breaking
breakatwhitespace=false,        % sets if automatic breaks should only happen at whitespace
%caption={\texttt{\lstname}},    % show the filename of files included with \lstinputlisting;
}

% New commands and enviromments
\newcommand{\TikZ}{Ti\emph{k}Z }
\newcommand{\PGF}{\textsc{PGF} }
\newcommand{\lcode}[1]{\texttt{#1}}  % Code in the same line.
\lstnewenvironment{code}{}{}  % Code in new line.
\newcommand{\fcode}[1]{\lstinputlisting[firstline=18]{#1}}  % Code in a file.
\newcommand{\example}[1]{  % For examples.
\begin{minipage}[htb]{0.55\textwidth}
    \fcode{#1}
\end{minipage} \hfill \vrule \hfill
\begin{minipage}[htb]{0.35\textwidth}
    \include{#1}
\end{minipage}
}


\begin{document}
\title[Introdu\c{c}\~{a}o ao LaTeX - 03/04]{Cursos CAMECC \\ Introdu\c{c}\~{a}o ao LaTeX para o Curso 29 \\ Licenciatura em Matem\'{a}tica \\ Aula 03}
\author{Raniere Silva}
\institute[CAMECC]{Centro Acad\^{e}mico dos Estudantes do IMECC}
\begin{frame}
    \titlepage
\end{frame}

\begin{frame}
    \frametitle{Licen\c{c}a}
    Este trabalho foi licenciado com a Licen\c{c}a Creative Commons Atribui\c{c}\~{a}o - CompartilhaIgual 3.0 N\~{a}o Adaptada. Para ver uma c\'{o}pia desta licen\c{c}a, visite \url{http://creativecommons.org/licenses/by-sa/3.0/} ou envie um pedido por carta para Creative Commons, 444 Castro Street, Suite 900, Mountain View, California, 94041, USA.
    %
    %Com base na obra dispon\'{i}vel em \url{https://github.com/r-gaia-cs/latex_with_vim/}. 
    %
    %Podem estar dispon\'{i}veis permiss\~{o}es adicionais ao \^{a}mbito desta licen\c{c}a em \url{https://github.com/r-gaia-cs/latex_with_vim/}.
    %
    \begin{center}
        \includegraphics[keepaspectratio=true]{../../figures/cc-by-sa.png}
    \end{center}
\end{frame}

\begin{frame}
    \frametitle{Sum\'{a}rio}
    \tableofcontents
\end{frame}

\section{\flang{Inline}}
\begin{frame}[fragile]
    \frametitle{\flang{Inline}}
    Exemplo: $1 + 1 = 2$.
    \pause

    \begin{code}
\documentclass{article}
\begin{document}
Exemplo: $1 + 1 = 2$.
\end{document}
    \end{code}
\end{frame}

\section{\flang{Displayed}}
\begin{frame}[fragile]
    \frametitle{\flang{Displayed}}
    Exemplo: $$1 + 1 = 2.$$
    \pause
    \begin{code}
\documentclass{article}
\begin{document}
Exemplo: $$1 + 1 = 2.$$
\end{document}
    \end{code}
\end{frame}

\section{Opera\c{c}\~{o}es b\'{a}sicas}
\begin{frame}[fragile]
    \frametitle{Opera\c{c}\~{o}es b\'{a}sicas} 
    \begin{code}
\documentclass{article}
\begin{document}
$1 \times 1 = 1$ \\
$1 \cdot 1 = 1$ \\
$1 / 1 = 1$ \\
$\frac{1}{1} = 1$
\end{document}
    \end{code}
\end{frame}

\begin{frame}[fragile]
    \frametitle{\'{I}ndices e expoentes} 
    \begin{code}
\documentclass{article}
\begin{document}
$a_11 \neq a_{11}$ \\
$1^0 = 1$
\end{document}
    \end{code}
\end{frame}

\begin{frame}[fragile]
    \frametitle{Delimitadores} 
    \begin{code}
\documentclass{article}
\begin{document}
$$\left( \frac{1}{1} \right) = 1$$
\end{document}
    \end{code}
\end{frame}

\begin{frame}[fragile]
    \frametitle{Matrizes} 
    \begin{code}
\documentclass{article}
\usepackage{amsmath}
\begin{document}
$$\begin{matrix}
    a_{11} & a_{12} \\
    a_{21} & a_{22}
\end{matrix}$$
\end{document}
    \end{code}
\end{frame}

\section{Equa\c{c}\~{o}es numeradas}
\begin{frame}[fragile]
    \frametitle{Equa\c{c}\~{o}es numeradas (1)}
    \begin{code}
\documentclass{article}
\usepackage{amsmath}
\begin{document}
\begin{align}
    1 + 1 &= 2
    \label{eq:test}
\end{align}
Na equa\c{c}\~{a}o \eqref{eq:test} \ldots
\end{document}
    \end{code}
\end{frame}

\begin{frame}[fragile]
    \frametitle{Equa\c{c}\~{o}es numeradas (2)}
    \begin{code}
\documentclass{article}
\usepackage{amsmath}
\begin{document}
\begin{align}
    1 + 1 &= 2 \tag{CAMECC}
    \label{eq:test}
\end{align}
Na equa\c{c}\~{a}o \eqref{eq:test} \ldots
\end{document}
    \end{code}
\end{frame}

\begin{frame}[fragile]
    \frametitle{Alinhamento de equa\c{c}\~{o}es}
    \begin{code}
\documentclass{article}
\usepackage{amsmath}
\begin{document}
\begin{align*}
    x + 1 &= 2 \\
    x + 1 - 1 &= 2 - 1 \\
    x &= 1
\end{align*}
% No lugar de \eqref{eq:test} ir\'{a} aparecer (??).
Na equa\c{c}\~{a}o \eqref{eq:test} \ldots
\end{document}
    \end{code}
\end{frame}

\begin{frame}
    \frametitle{Fim}
    \begin{center}
        Perguntas.
    \end{center}
\end{frame}
\end{document}
